% ============================================================
% Chapter 1: Introduction
%   Section: Motivation
% Included from chapters/01-introduction/introduction.tex
% Source: docs/idea-de-negocio.md (principles), docs/roadmap.md
% (the thesis in one sentence), docs/decisions/deployment.md (cost).
% ============================================================

\section{Motivation}
\label{sec:motivation}

\paragraph{Verification inside the flow, not after it.} A fact-check
usually arrives after a story has circulated, and is read by fewer
people than the story it corrects. A publication of constructive news
has a stronger reason than most to check before it publishes: its
value to readers is that the good news it reports is true, and a
single inflated figure or misattributed result costs it that value.
The motivation of this work is therefore to make verification a step
of the flow that finds and prepares the news, so that every published
article carries a verdict per claim, with the evidence behind it,
rather than an editorial label nobody can inspect. The same flow also
does the work before verification that takes a journalist's time ---
reading many sources to find the few stories worth offering --- so
that the two can be measured together.

\paragraph{Failing towards caution.} An automated verifier that is
sometimes confidently wrong is worse than none: a reader cannot tell
its right answers from its wrong ones. The system is designed to
prefer ``not verified'' to a confident verdict it cannot ground, to
report a search or model failure as a failure rather than as a
verdict, and to show the evidence and reasoning behind every answer.

\paragraph{Low cost and the open web.} The evidence for a claim about
last week's study or this month's project is not in any fixed corpus:
it is on the web, if anywhere, on the day the claim is checked. The
system therefore searches the open web at verification time, as a
person would. It does so without paid services: the search goes
through a self-hosted metasearch engine, and the language model,
entity recogniser, sentiment classifier and embedding model all run
locally from open weights. The whole stack runs on a laptop, or on a
single cloud virtual machine of about 115~USD a month
(Section~\ref{sec:dev-deployment}). A small newsroom or a single
researcher can run it and reproduce every number in this thesis, and
no result depends on a provider's undisclosed model or index. The
price is a small model and a search engine that is sometimes refused
by the engines behind it, and both are measured rather than assumed.
